\chapter{Reset architecture}
\label{ch:reset}

\section{The rule}

One reset signal, \sig{rst\_n\_i}, asynchronous and active low, reaches every
sequential element in the delivery. The policy applied to it is deliberately
absolute:

\begin{rtlnote}
\textbf{Every flop in the design is reset, and no flop leaves reset holding X.}
This includes state that a correctness argument alone would not require to be
reset --- pipeline-register payloads covered by a valid bit, branch-predictor
arrays covered by a valid bit, the AXI command latches, and the interrupt
controller's nesting-stack arrays.
\end{rtlnote}

\section{Why the rule is absolute}

The narrower rule --- reset only what correctness requires --- is a defensible
one, and it is what an earlier revision of this design followed. It was
tightened for four reasons, in descending order of how much they matter.

\paragraph{1. An unreset flop can reach a bus output.} This is not
hypothetical; it was present in the delivered design and is the finding that
motivated the change. In the shared AXI4-Lite slave, \sig{BRESP} and
\sig{RRESP} are combinational functions of a stored error flag, and they are
driven \emph{whether or not} the corresponding \sig{VALID} is asserted. With
those flags unreset, the slave presented X on its response codes from time zero
until its first transaction completed. A protocol monitor cannot distinguish
that from a real violation, and neither can a reviewer reading a waveform.

\paragraph{2. Two-state and four-state simulation must agree.} A design with
unreset state behaves differently under a four-state simulator, which propagates
X, and a two-state simulator or synthesis-derived netlist, which does not. Any
argument of the form ``the valid bit prevents the X from mattering'' has to hold
for every consumer of every unreset bit, and it has to keep holding after the
next change. Removing the X removes the argument.

\paragraph{3. Waveforms are evidence.} Verification here rests on reading
waveforms and on protocol monitors that sample real signals. A fetch stage whose
program counter reads X until the first instruction arrives is an unreadable
waveform in exactly the window --- the first cycles after release --- where
reset behaviour is being reviewed.

\paragraph{4. Any register may become architecturally visible.} The interrupt
controller's nesting stack was internal until \reg{NEST\_STATUS.TOP\_ID} and
\reg{ACTIVE\_VEC.ID} were added, at which point unreset entries became a value
software could read. The rule removes the need to re-audit that question every
time a status register grows a field.

\paragraph{What it costs.} The cost is real and is stated rather than hidden:
the branch predictor's payload arrays and the return-address stack move onto the
reset tree, roughly 7.2\,kbit at the default configuration of 128 entries and a
depth-8 stack. No synthesis was run for this delivery, so no area or timing
figure is quoted. If a future target cannot afford it, the payload
\texttt{always} blocks in \file{branch\_predictor.v} are the only places to
change, and the design remains functionally identical --- the predictor is never
believed before its valid bit is set.

\paragraph{Packed, not arrays.} The branch predictor's five state elements are
declared as packed vectors rather than unpacked arrays, and are addressed with
an indexed part-select. That is the same hardware, but it lets each one reset in
a \emph{single} assignment instead of a loop over 128 entries. The distinction
is not cosmetic: a non-blocking assignment to an unpacked array inside a
\texttt{for} loop is only accepted by a tool that can unroll that loop, and the
unroll budget is a version-dependent default --- 128 entries is past the value
older Verilator releases use. Writing the reset without a loop removes the
dependency altogether. The local Verilator flow additionally pins
\texttt{-{}-unroll-count 64} so that a future loop which would exceed the
budget fails on the developer's machine rather than in CI.

\section{One block, one owner}

A second rule sits alongside the first and is what makes the reset values
readable at all:

\begin{rtlnote}
\textbf{Exactly one \texttt{always} block drives each register or memory.} A
register's reset value and every rule that updates it are therefore always found
together, in one place, and no register can be assigned from two processes.
\end{rtlnote}

\section{Reset values, block by block}

The tables below list every reset value in the delivery, read from the source.
Where the value is not zero it is because zero would be the wrong behaviour, not
a neutral one; the reason is given in each case.

{\small
\begin{longtable}{@{}L{0.20\linewidth}llL{0.40\linewidth}@{}}
\caption{CPU core reset values (\file{hdl/cpu\_top.v}, \file{fetch\_unit.v}, \file{branch\_predictor.v}, \file{regfile.v}, \file{lsu.v})}
\label{tab:cpureset}\\
\toprule
\textbf{Register} & \textbf{Width} & \textbf{Reset} & \textbf{Note} \\
\midrule
\endfirsthead
\toprule
\textbf{Register} & \textbf{Width} & \textbf{Reset} & \textbf{Note} \\
\midrule
\endhead
\bottomrule
\endfoot
\bottomrule
\endlastfoot
\multicolumn{4}{@{}l@{}}{\emph{IF/DX pipeline register}} \\
\sig{ifdx\_valid\_q}   & 1  & 0 & the bubble bit \\
\sig{ifdx\_instr\_q}   & 32 & \texttt{0x0000\_0013} & the canonical NOP; the
  decoder is combinational, so this keeps \sig{MemRead}, \sig{MemWrite} and
  \sig{illegal} defined from time zero \\
\sig{ifdx\_pc\_q}      & 32 & \sig{RESET\_PC} & feeds the PC+4 adder, the branch
  target adder and the predictor update port combinationally \\
\sig{ifdx\_ptk\_q}     & 1  & 0 & no prediction \\
\sig{ifdx\_ptg\_q}     & 32 & \sig{RESET\_PC} & \\
\sig{ifdx\_fault\_q}   & 1  & 0 & \\
\midrule
\multicolumn{4}{@{}l@{}}{\emph{DX/WB pipeline register}} \\
\sig{dxwb\_valid\_q}   & 1  & 0 & \\
\sig{dxwb\_rd\_q}      & 5  & 0 & x0, harmless as a writeback target \\
\sig{dxwb\_regwrite\_q}& 1  & 0 & \\
\sig{dxwb\_wbsel\_q}   & 2  & \sig{WB\_ALU} & \\
\sig{dxwb\_result\_q}  & 32 & 0 & feeds \sig{wb\_data} and the S3$\rightarrow$S2
  forwarding path combinationally \\
\sig{dxwb\_mem\_q}     & 32 & 0 & \\
\sig{dxwb\_pc4\_q}     & 32 & 0 & \\
\midrule
\multicolumn{4}{@{}l@{}}{\emph{Trap and interrupt handshake}} \\
\sig{in\_irq}          & 1 & 0 & \\
\sig{cpu\_irq\_ack\_o}  & 1 & 0 & output; must be low out of reset \\
\sig{cpu\_irq\_eoi\_o}  & 1 & 0 & output; must be low out of reset \\
\sig{cpu\_in\_trap\_o}  & 1 & 0 & output \\
\midrule
\multicolumn{4}{@{}l@{}}{\emph{Fetch stage}} \\
\sig{ar\_pending\_q}   & 1  & 0 & \\
\sig{inflight\_q}      & 1  & 0 & \\
\sig{discard\_q}       & 1  & 0 & \\
\sig{issued\_pc\_q}    & 32 & \sig{RESET\_PC} & drives \sig{ARADDR} \\
\sig{hold\_valid\_q}   & 1  & 0 & \\
\sig{hold\_instr\_q}   & 32 & \texttt{0x0000\_0013} & the same NOP, so even an
  invalid holding slot decodes harmlessly \\
\sig{hold\_pc\_q}      & 32 & \sig{RESET\_PC} & \\
\sig{hold\_fault\_q}   & 1  & 0 & \\
\sig{npc\_q}           & 32 & \sig{RESET\_PC} & the deferred fetch address \\
\midrule
\multicolumn{4}{@{}l@{}}{\emph{Branch predictor} (\sig{ENTRIES}~=~128,
  \sig{RAS\_DEPTH}~=~8)} \\
\sig{valid}  & 128 & 0 & one bit per entry \\
\sig{isret}  & 128 & 0 & one bit per entry \\
\sig{state}  & 128 & 0 & the 1-bit saturating direction, one bit per entry \\
\sig{tag}    & 128\,$\times$\,23 & 0 & payload, reset for X-freedom \\
\sig{target} & 128\,$\times$\,32 & 0 & payload; read combinationally on every
  lookup, hit or miss \\
\sig{ras[]}  & 8\,$\times$\,32   & 0 & payload \\
\sig{sp\_q}   & 3 & 0 & \\
\sig{cnt\_q}  & 4 & 0 & stack empty \\
\midrule
\multicolumn{4}{@{}l@{}}{\emph{Register file and load/store unit}} \\
\sig{regs[]}    & 32\,$\times$\,32 & 0 & REQ10 \\
\sig{state\_q}  & 2  & \sig{S\_IDLE} & \\
\sig{ar\_sent\_q}, \sig{aw\_sent\_q}, \sig{w\_sent\_q} & 1 each & 0 & \\
\sig{addr\_q}   & 32 & 0 & command latch; drives \sig{AWADDR}/\sig{ARADDR} \\
\sig{wdata\_q}  & 32 & 0 & command latch; drives \sig{WDATA} \\
\sig{wstrb\_q}  & 4  & 0 & command latch; drives \sig{WSTRB} \\
\end{longtable}}

{\small
\begin{longtable}{@{}L{0.26\linewidth}lL{0.42\linewidth}@{}}
\caption{CSR file reset values (\file{hdl/csr\_file.v})}
\label{tab:csrreset}\\
\toprule
\textbf{CSR} & \textbf{Reset} & \textbf{Note} \\
\midrule
\endfirsthead
\toprule
\textbf{CSR} & \textbf{Reset} & \textbf{Note} \\
\midrule
\endhead
\bottomrule
\endfoot
\bottomrule
\endlastfoot
\reg{mstatus}   & \texttt{0x0000\_1800} & MIE = 0 and MPIE = 0, so interrupts are
  off at boot; MPP is hardwired to \texttt{11} (M-mode) and is what makes the
  reset value non-zero \\
\reg{mie}       & \texttt{0x0000\_0000} & no source enabled \\
\reg{mtvec}     & \texttt{0x0000\_0000} & direct mode, base 0 \\
\reg{mscratch}  & \texttt{0x0000\_0000} & \\
\reg{mepc}      & \texttt{0x0000\_0000} & \\
\reg{mcause}    & \texttt{0x0000\_0000} & \\
\reg{mtval}     & \texttt{0x0000\_0000} & \\
\reg{mcycle}, \reg{minstret}, \reg{mhpmcounter3..7} and their upper halves
                & 0 (64-bit) & \reg{mcycle} is free-running, so it is already
  non-zero by the first access \\
\midrule
\reg{misa}      & \texttt{0x4000\_0100} & hardwired, not a flop: MXL = 32,
  extension I \\
\reg{mhartid}   & \sig{HART\_ID} & hardwired from the parameter; unaffected by
  reset \\
\reg{mvendorid}, \reg{marchid}, \reg{mimpid} & 0 & hardwired \\
\reg{mip}       & driven & a function of the PIC lines, not stored \\
\end{longtable}}

{\small
\begin{longtable}{@{}L{0.19\linewidth}llL{0.38\linewidth}@{}}
\caption{Interrupt controller reset values (\file{hdl/pic.v})}
\label{tab:picreset}\\
\toprule
\textbf{State} & \textbf{Width} & \textbf{Reset} & \textbf{Note} \\
\midrule
\endfirsthead
\toprule
\textbf{State} & \textbf{Width} & \textbf{Reset} & \textbf{Note} \\
\midrule
\endhead
\bottomrule
\endfoot
\bottomrule
\endlastfoot
\sig{src\_config[]}  & 16\,$\times$\,32 & 0 & level-triggered, band 0,
  intra-priority 0, deadline disabled \\
\sig{sw\_pend}       & 16 & 0 & \\
\sig{int\_enable}    & 16 & 0 & \textbf{the disarming condition:} nothing can
  reach the CPU until software opts a source in \\
\sig{band\_cfg}      & 8  & \texttt{0x1B} & band0 = 3 (most urgent) down to
  band3 = 0, the natural ordering. All-zero would make the four bands equally
  urgent, which is not a neutral default but a broken one \\
\sig{nest\_max}      & 5  & 8 & a limit of 0 would mask every offer for ever;
  writes are additionally clamped to $[1,16]$ \\
\sig{esc\_cfg}       & 9  & 0 & jump mode, target band 0, single escalation \\
\sig{spurious\_log}  & 16 & 0 & \\
\sig{int\_status}    & 3  & 0 & \\
\sig{irq\_src\_q}    & 16 & 0 & edge-detect sample \\
\sig{edge\_pend}     & 16 & 0 & \\
\sig{active}         & 16 & 0 & \\
\sig{escalated}      & 16 & 0 & \\
\sig{spurious}       & 16 & 0 & \\
\sig{eff\_band[]}    & 16\,$\times$\,2  & 0 & \\
\sig{ddl\_cnt[]}     & 16\,$\times$\,16 & 0 & \\
\sig{stack\_id[]}    & 16\,$\times$\,4  & 0 & read by \reg{NEST\_STATUS.TOP\_ID}
  and \reg{ACTIVE\_VEC.ID} \\
\sig{stack\_key[]}   & 16\,$\times$\,10 & 0 & the preemption threshold \\
\sig{depth}          & 5 & 0 & \\
\sig{cpu\_irq\_o}     & 1 & 0 & output \\
\sig{cpu\_irq\_vec\_o} & 4 & 0 & output \\
\sig{cpu\_irq\_vec\_d} & 4 & 0 & \\
\end{longtable}}

\begin{table}[H]
\centering
\caption{Machine timer and shared slave reset values
(\file{hdl/mtimer.v}, \file{hdl/axi\_lite\_slave.v})}
\label{tab:tmrreset}
\begin{tabular}{@{}lllL{0.34\linewidth}@{}}
\toprule
\textbf{State} & \textbf{Width} & \textbf{Reset} & \textbf{Note} \\
\midrule
\sig{mtime\_q}     & 64 & 0 & free-running from the first clock after release \\
\sig{mtimecmp\_q}  & 64 & all ones & \textbf{the timer is born disarmed}: the
  counter cannot reach the compare value, so the interrupt line cannot fire
  before software arms it. A reset of 0 would fire immediately \\
\sig{irq\_o}        & 1  & 0 & output \\
\midrule
\sig{aw\_got\_q}, \sig{w\_got\_q} & 1 each & 0 & \\
\sig{awoff\_q}     & 6  & 0 & command payload \\
\sig{wdata\_q}     & 32 & 0 & command payload \\
\sig{wstrb\_q}     & 4  & 0 & command payload \\
\sig{s\_axi\_bvalid\_o}, \sig{s\_axi\_rvalid\_o} & 1 each & 0 & required low during
  reset by IHI0022E A3.1.2 \\
\sig{bresp\_err\_q}& 1  & 0 & \textbf{load-bearing:} \sig{BRESP} is driven from
  this flop whether or not \sig{BVALID} is high \\
\sig{rresp\_err\_q}& 1  & 0 & same argument for \sig{RRESP} \\
\sig{s\_axi\_rdata\_o}& 32 & 0 & \\
\bottomrule
\end{tabular}
\end{table}

\section{Timing of assertion and release}

The reset is asynchronous, so assertion takes effect at the instant it occurs
rather than at the next clock edge; release takes effect on the following rising
edge. The clock and reset generator used by the benches releases reset at
123\,ns with a 10\,ns clock period --- deliberately between clock edges, so that
a design accidentally relying on a synchronous release would be exposed.

One output is gated combinationally rather than only registered:
\sig{ibus\_arvalid\_o} is qualified with \sig{rst\_n\_i}, because the fetch-issue
decision is combinational and would otherwise request the reset vector while
reset was still asserted. AXI4-Lite forbids \sig{VALID} during reset, and a
slave holding \sig{READY} high in reset would have accepted a phantom address.
This was a real defect, found by the assertion layer on its first run.

\section{How the rule is verified}

The reset policy is not asserted in this document; it is measured. The evidence
is in Chapter~\ref{ch:verifplan}, sections \ref{sec:vp-picreset},
\ref{sec:vp-mtimer} and \ref{sec:vp-csr}. In outline, three properties are
checked for each block:

\begin{enumerate}[leftmargin=*,itemsep=2pt]
\item every mapped register reads back the value listed in the tables above, on
its first access after release;
\item nothing in the block --- including state with no register view, probed
directly through hierarchical references --- holds X after release;
\item a reset asserted \emph{off a clock edge} from a fully configured, actively
nested state returns every one of those registers to the same value, which is
what distinguishes a real asynchronous reset from a simulation initial value.
\end{enumerate}

The check has been confirmed to have teeth: run against the previous revision of
the RTL, it fails on precisely the flops that revision left unreset --- the two
AXI response-code flags and the interrupt controller's nesting-stack arrays.
